\documentclass[10pt,a4paper]{amsart}
\usepackage[margin=19mm]{geometry}
\usepackage{amsmath,amssymb,mathtools}
\usepackage[scr=rsfs,cal=euler]{mathalfa}
\usepackage{xfrac}
\usepackage[unicode,hidelinks,bookmarks=false]{hyperref}
\newcommand{\ie}{i.e.\ }
\newcommand{\cf}{cf.\ }
\newcommand{\resp}{respectively\ }
\newcommand{\Subsection}{\subsection}
\providecommand{\nobreakdash}{\nobreak\hbox{-}}
\setlength{\emergencystretch}{2em}
\multlinegap0pt
\usepackage{bm}
\usepackage{bbm}
\usepackage{dsfont}
\usepackage{xspace}
\usepackage{cancel}
\usepackage{mathtools}
\usepackage{amsthm}
\makeatletter
\@ifundefined{equation}{\newtheorem{equation}{Equation}}{}
\@ifundefined{Lemma}{\newtheorem{Lemma}[equation]{Lemma}}{}
\@ifundefined{remark}{\newtheorem{remark}[equation]{Remark}}{}
\makeatother
\DeclarePairedDelimiter{\abs}{\lvert}{\rvert}% absolute value
\newcommand{\Cost}{\text{Cost}}
\title[On The Cost Function Associated With Legendrian Knots]{On The Cost Function Associated With Legendrian Knots}
\author{Dheeraj Kulkarni, Tanushree Shah, Monika Yadav}
\date{Source: arXiv:2307.13963v3 (CC BY 4.0)}
\begin{document}
\maketitle

\noindent\emph{Source and licence.} On The Cost Function Associated With Legendrian Knots. Version arXiv:2307.13963v3 (CC BY 4.0), distributed under the Creative Commons Attribution 4.0 International licence, as recorded in the arXiv licence metadata for this exact version and shown on the abstract page https://arxiv.org/abs/2307.13963v3. Full rights notice: licenses/arxiv-2307.13963-CC-BY-4.0.txt.\par\smallskip

\noindent\textit{Editorial scope.} Counted unit: the Lemma whose source label is \texttt{CostL21} in the article (source0.tex lines 769--771, source label \texttt{CostL21}), delivered together with the article's own complete original proof of it (source0.tex lines 772--839, one \texttt{proof} environment of 68 lines). No mathematics is written, repaired or re-derived: statement and proof are the article's own text line for line. Required context retained uncounted: CostL8 (lines 491--493); NewRemark2 (lines 498--501). Every definition, display and label of those lines is retained unchanged. Omitted from this package and not redistributed: 1--490, 494--497, 502--768, 840--962. Nothing inside a retained excerpt is omitted or altered: every retained body is delivered exactly as the article's lines are written, so no omission receipt of any kind accompanies this package. Citations: the retained excerpts carry no citation command at all, so no bibliography entry of the article is transcribed and no reference is redistributed. Reference adaptation: none was needed and none was made; every reference inside the delivered unit resolves inside it, so no unresolved reference or citation is produced. Printed slips kept literal: no printed word, symbol, hypothesis or number of the counted statement or of its proof was altered. Layout and class: the article's own class is \texttt{amsart} with packages including inputenc, fontenc, lmodern, amsmath, amsthm, amssymb, amsfonts, latexsym, amsxtra, xy, nicefrac, mathtools, enumitem, microtype; the wrapper is the lane's pinned \texttt{amsart} editorial wrapper at 10pt on A4, which restates the theorem-like environments the retained text needs and adds the article's own macro definitions that the retained text needs (\texttt{\textbackslash Cost}, \texttt{\textbackslash abs}), with extra wrapper packages (bm, bbm, dsfont, xspace, cancel, mathtools). The delivered numbering is the unit's own. Assets: no retained span contains a figure environment, a graphics-inclusion command, a PDF-page inclusion, a table or a TikZ picture; no image file is copied into this package, the package declares no asset dependency and no image file is redistributed with it. Licence: version arXiv:2307.13963v3 of this article is distributed under the Creative Commons Attribution 4.0 International licence, as recorded in the arXiv licence metadata for this exact version and shown on the abstract page https://arxiv.org/abs/2307.13963v3. A full rights notice is delivered at licenses/arxiv-2307.13963-CC-BY-4.0.txt. Characters: every retained excerpt is pure ASCII, so the delivered engine, XeLaTeX with its default Latin Modern text font, reports no missing character.
\begin{Lemma}\label{CostL8}
    Let $K$ and $\Tilde{K}$ be Legendrian knots of the same topological knot type. Then $$\Cost(K,\Tilde{K})\geq \max\{\abs{tb(K)-tb(\Tilde{K})},\abs{rot(K)-rot(\Tilde{K})}\}.$$
\end{Lemma}

\begin{remark}\label{NewRemark2}
 Note that if $K=(S^+)^m(S^-)^n(L)$ then $\Cost(K,L)=m+n$, where $m$ and $n$ are positive integers. For the justification, we clearly have $
   \Cost(K,L)\leq m+n$ and $tb(K)=tb(L)-m-n$. Now, by Lemma \ref{CostL8}, we have $\Cost(K,L)\geq |tb(K)-tb(L)|= m+n$.
\end{remark}

\begin{Lemma}\label{CostL21}
Let $K_1,L_1$ and $K_2,L_2$ be Legendrian knots of two different topological knot types such that $K_1=(S^+)^p(S^-)^n(L_1)$ and $ L_2=(S^+)^q(S^-)^{m}(K_2)$. Then $$\Cost(K_1\#K_2,L_1\#L_2)=|p-q|+|n-m|.$$
\end{Lemma}

\begin{proof}
We will prove this Lemma in cases. 
    
    \noindent \textbf{Case I:} $p\geq q,n\geq m$.
    
  \noindent We prove $K_1\#K_2$ destabilizes to $L_1\#L_2$ as follows:
   \begin{align*}
       K_1\#K_2&=(S^+)^p(S^-)^n(L_1\#K_2)\\
       &=(S^+)^{p-q}(S^-)^{n-m}(L_1\#(S^+)^q(S^-)^{m}(K_2))\\
       &=(S^+)^{p-q}(S^-)^{n-m}(L_1\#L_2).&&
   \end{align*} 
   Using Remark \ref{NewRemark2}, we have
   \begin{align*}
       \Cost(K_1\#K_2,L_1\#L_2)&=p-q+n-m\\
       &=|\Cost(K_1,L_1)-\Cost(K_2,L_2)|.
   \end{align*}
   
 \noindent \textbf{Case II:} $p\leq q,n\leq m$.

   \noindent Using the similar technique as used in Case I, we prove that $K_1\#K_2$ destabilizes to $L_1\#L_2$.
   \begin{align*}
       L_1\#L_2&=(S^+)^q(S^-)^{m}(L_1\#K_2)\\
       &=(S^+)^{q-p}(S^-)^{m-n}((S^+)^p(S^-)^n(L_1)\#K_2)\\
       &=(S^+)^{q-p}(S^-)^{m-n}(K_1\#K_2).&&
   \end{align*} 
   Therefore, by Remark \ref{NewRemark2}
   \begin{align*}
       \Cost(K_1\#K_2,L_1\#L_2)&=q-p+m-n\\
       &=|\Cost(K_1,L_1)-\Cost(K_2,L_2)|.
   \end{align*}

   \noindent \textbf{Case III:} $p\geq q,n\leq m$.

\noindent Since $ rot(K_1)=rot(L_1)-n+p$, and $,rot(L_2)=rot(K_2)-m+q$, using Lemma \ref{CostL8}, we get
\begin{align*}
    \Cost(K_1\#K_2,L_1\#L_2)&\geq |rot(K_1\#K_2)-rot(L_1\#L_2)|\\
    &=p-q+m-n.
\end{align*}
The reverse inequality follows from the computations done below. 

   \begin{align*}
       (S^-)^{m-n}(K_1\#K_2)&=(S^-)^{m-n}((S^+)^p(S^-)^n(L_1)\#K_2)\\
       &=(S^+)^p(S^-)^{m}(L_1\#K_2)\\
       &=(S^+)^{p-q}(L_1\#(S^+)^q(S^-)^m(K_2))\\
       &=(S^+)^{p-q}(L_1\#L_2).
   \end{align*}

   
   \noindent \textbf{Case IV:} $p\leq q,m\leq n$.
   

      Using Lemma \ref{CostL8}, we have
\begin{align*}
    \Cost(K_1\#K_2,L_1\#L_2)&\geq |rot(K_1\#K_2)-rot(L_1\#L_2)|\\
    &=|p-q+m-n|=q-p+n-m.
\end{align*}

   The reverse inequality follows from the following computations:
\begin{align*}
       (S^+)^{q-p}(K_1\#K_2)&=(S^+)^{q-p}((S^+)^p(S^-)^n(L_1)\#K_2)\\
       &=(S^+)^q(S^-)^{n}(L_1\#K_2)\\
       &=(S^-)^{n-m}(L_1\#(S^+)^q(S^-)^m(K_2))\\
       &=(S^-)^{n-m}(L_1\#L_2).
   \end{align*}
  

   Hence, we conclude from the above cases that $$\Cost(K_1\#K_2,L_1\#L_2)=|p-q|+|n-m|.$$
\end{proof}

\end{document}
